\documentclass[a4paper,11pt]{article}

\usepackage[utf8]{inputenc}
\usepackage[T1]{fontenc}
\usepackage[french]{babel}
\usepackage{lmodern}
\usepackage{microtype}
\usepackage{amsmath}
\usepackage{booktabs}
\usepackage{longtable}
\usepackage{array}
\usepackage{geometry}
\usepackage{xcolor}
\usepackage{hyperref}

\geometry{margin=2.5cm}
\hypersetup{
  colorlinks=true,
  linkcolor=blue!45!black,
  urlcolor=blue!45!black,
  citecolor=blue!45!black,
  pdftitle={ASTRA P78 Results},
  pdfauthor={ASTRA / Representations Procedurales Adressables}
}

\setlength{\parindent}{0pt}
\setlength{\parskip}{0.62em}
\renewcommand{\arraystretch}{1.12}
\setlength{\emergencystretch}{2em}

\newcommand{\Metric}[2]{#1 & #2\\}
\newcommand{\Decision}{\texttt{RECALIBRATE\_P78\_LEVEL1\_TOPOLOGY}}

\title{\textbf{ASTRA-P78 Results}\\Level-1 Virtual Space and Universal Fiber Store}
\author{ASTRA / Repr\'esentations Proc\'edurales Adressables}
\date{2026-05-01}

\begin{document}
\maketitle

\begin{abstract}
P78 teste la topologie de niveau 1 de l'espace d'adressage virtuel. Le meilleur
candidat observe est l'espace hybride multi-index: il ameliore le ratio living
tout en gardant un store universel avec fallback brut explicite, guard refuse,
reopen equivalence et aucun drift dur. La decision reste prudente car le meilleur
cout d'adressage revient encore au trie de chemins.
\end{abstract}

\section{Objectif}

P78 deplace l'analyse au-dessus des fibres: il compare la topologie de l'espace
d'adressage virtuel de niveau 1. L'espace virtuel ne doit pas etre materialise
globalement. Il est calcule localement quand une adresse est atteinte.

Le jalon ajoute aussi un store universel capable d'accepter des fichiers texte,
code, JSON, CSV, logs, binaires structures, binaires inconnus via fallback brut,
et guard incompressible refuse.

\section{Topologies}

Les topologies testees sont:
\begin{itemize}
  \item grille 2D et grille 3D;
  \item arbre hierarchique;
  \item trie de chemins;
  \item DAG adresse par contenu;
  \item graphe d'adresses;
  \item espace produit type;
  \item espace hybride multi-index.
\end{itemize}

Le candidat gagnant par ratio est l'espace hybride multi-index. Le meilleur cout
d'adressage est le trie de chemins. Cette dissociation motive la recalibration
plutot qu'une promotion.

\section{Campagnes living-memory}

Les campagnes significatives utilisent `10,485,760` octets de donnees sources,
un store vivant, open/read/query/update/delete/audit/compact/close/reopen, et
verification du guard incompressible.

\begin{longtable}{@{}p{0.58\linewidth}r@{}}
\toprule
Metrique & Valeur\\
\midrule
\Metric{target source bytes}{10,485,760}
\Metric{actual source bytes}{10,485,760}
\Metric{cycles standard / ambitious}{10 / 25}
\Metric{queries standard / ambitious}{10,000 / 50,000}
\Metric{updates standard / ambitious}{1,000 / 5,000}
\Metric{deletes standard / ambitious}{100 / 500}
\Metric{best level-1 topology}{hybrid multi-index}
\Metric{phase green / yellow / red}{4 / 2 / 2}
\bottomrule
\end{longtable}

\section{Espace virtuel}

Les octets virtuels ci-dessous sont des equivalents de materialisation, pas des
octets stockes.

\begin{longtable}{@{}p{0.58\linewidth}r@{}}
\toprule
Metrique & Valeur\\
\midrule
\Metric{level1 declared addresses}{281,600,000,000}
\Metric{level1 reachable addresses}{247,808,000,000}
\Metric{level1 effective addresses}{222,464,000,000}
\Metric{virtual cells}{10,000}
\Metric{virtual fibers}{40,000}
\Metric{virtual chunks}{40,000}
\Metric{virtual versions}{4}
\Metric{virtual effective bytes equivalent}{96,415,629}
\Metric{materialization avoidance ratio}{11.639128}
\Metric{limiting factor}{index size}
\bottomrule
\end{longtable}

\section{Resultats}

\begin{longtable}{@{}p{0.58\linewidth}r@{}}
\toprule
Metrique & Valeur\\
\midrule
\Metric{cold persisted bytes}{1,268,502}
\Metric{runtime peak bytes}{559,633}
\Metric{ratio living}{5.340001}
\Metric{address lookup p95 steps}{8.000}
\Metric{address lookup p95 bytes read}{4,608}
\Metric{CRUD success rate}{1.000000}
\Metric{roundtrip / decode / retrieval}{1.000000}
\Metric{raw fallback bytes}{327,680}
\Metric{guard source bytes}{524,288}
\Metric{guard store bytes}{0}
\Metric{drift status}{NO\_DRIFT}
\Metric{reopen equivalence}{true}
\bottomrule
\end{longtable}

\section{Universal Codec}

Le store accepte les familles de fichiers via codecs specialises ou fallback
brut explicite: grammar tokens, JSON path, CSV sparse, dictionnaire texte,
deduplication de chunks, DAG de contenu, fallback brut et guard refuse.

Le fallback brut est compte comme stockage reel. Il ne contribue pas a un faux
gain. Le guard incompressible reste:

\texttt{NO\_GO\_GUARD\_INCOMPRESSIBLE\_REFUSED}.

\section{Decision}

\Decision

Raison principale: l'espace hybride multi-index maximise \texttt{ratio\_living}, mais
le trie de chemins reste meilleur pour le cout d'adressage. P78 recommande donc
de tester en P79 un routeur de niveau 1 par type de fichier et feature locale.

\section{Limites}

Les corpus binaires image/video sont des substituts deterministes locaux, pas
des datasets externes lourds. La recherche niveau 1 reste coarse. Les octets
virtuels sont des equivalents de materialisation; ils ne doivent pas etre lus
comme de l'information stockee gratuitement.

\section{P79}

P79 devrait evaluer un routeur de niveau 1: trie pour chemins, DAG de contenu
pour chunks et deduplication, espace produit type pour namespace/type/objet,
graphe pour relations, et hybride multi-index seulement quand son overhead est
justifie.

\end{document}
